\color{unidarkgreen}{\subsection{Дистанционная защита (ДЗ)}}
\color{black}

\begin{enumerate}[label=\arabic{section}.\arabic{subsection}.\arabic{enumi}, labelsep=4pt, leftmargin=0pt, itemindent=57pt, itemsep=0pt, parsep=5pt] % если что можно крутить itemsep
%%%%%%%%%%%%%%%%%%%%%%%%%%%%%%%%%%%%%%%%%%%%%%%%%%%%%%%%%%%%%%%%%%%%%%  ОБЩИЕ СВЕДЕНИЯ %%%%%%%%%%%%%%%%%%%%%%%%%%%%%%%%%%%%%%%%%%%%%%%%%%
\item Общие сведения

\begin{enumerate}[label=\arabic{section}.\arabic{subsection}.\arabic{enumi}.\arabic*, labelsep=4pt, leftmargin=0em, itemindent=65pt, parsep=0pt]

\item
ИЭУ содержит шесть ступеней ДЗ от междуфазных замыканий и две ступени ДЗ от замыканий на
землю. Дистанционная  защита  (ДЗ)  предназначена  для  селективного  отключения  поврежденного  элемента
электрической  сети  при  всех видах  коротких  замыканий (КЗ),  реагируя  на снижение измеряемого
сопротивления в защищаемой зоне.

\item 
ДЗ селективно срабатывает:
\begin{itemize}
    \item при всех видах КЗ; 
    \item при переходе внешнего КЗ во внутреннее; 
    \item при КЗ, сопровождающемся снижением измеряемого напряжения до нуля.
\end{itemize}

\item 
Обеспечивается несрабатывание защиты:
\begin{itemize} 
\item при качаниях и асинхронном ходе; 
\item при неисправностях цепей напряжения;
\item при внешних КЗ с насыщением трансформаторов тока (при соблюдении заявленных производителем требований к ТТ; 
\item при бросках тока намагничивания без КЗ;
\item при внешних КЗ, даже если измеряемое напряжение снижается до нуля.
\end{itemize}

\item Каждая из ступеней ДЗ от междуфазных КЗ включает в себя три реле сопротивления (РС), использующие
для замера сопротивления  комплексные значения линейных напряжений и соответствующих им линейных
токов:
% ФОРМУЛЫ
\begin{align}
\vec Z_{AB} =  \frac{\vec U_A - \vec U_B}{\vec I_A - \vec I_B}\label{eq:groupZ1}, \\[10pt]
\vec Z_{BC} =  \frac{\vec U_B - \vec U_C}{\vec I_B - \vec I_C}\label{eq:groupZ2}, \\[10pt]
\vec Z_{CA} =  \frac{\vec U_C - \vec U_A}{\vec I_C - \vec I_A}\label{eq:groupZ3}, 
\end{align}
\begin{eqexpl}[20mm]
    \item{$ \vec U_{A}, \vec U_{B}, \vec U_{C} $} комплексные значения напряжений фаз А, В, С;
	\vspace{2pt}
    \item{$ \vec I_{A}, \vec I_{B}, \vec I_{C} $} комплексные значения токов фаз А, В, С.
\end{eqexpl}
\vspace{4pt}

\item 
Комплексные значения сопротивлений по контуру «фаза-земля» рассчитываются с учетом компенсации
тока нулевой последовательности:
% ФОРМУЛЫ
\begin{align}
\vec Z_{A} =  \frac{\vec U_A }{\vec I_A + \vec K \cdot  3 \vec I_0}\label{eq:groupZ11}, \\[10pt]
\vec Z_{B} =  \frac{\vec U_B }{\vec I_B + \vec K \cdot  3 \vec I_0}\label{eq:groupZ12}, \\[10pt]
\vec Z_{C} =  \frac{\vec U_C }{\vec I_C + \vec K \cdot  3 \vec I_0}\label{eq:groupZ13}, 
\end{align}
\begin{eqexpl}[5mm]
    \item{$ 3 \vec I_{0} $}  комплексное  значение  утроенного  тока  нулевой  последовательности  защищаемого объекта (расчетное значение);
    \item{$ \vec K $}  комплексный коэффициент компенсации тока нулевой последовательности защищаемого объекта.
\end{eqexpl}
\vspace{4pt}

\item 
В ИЭУ коэффициент  $\vec K$  задается двумя уставками, \textbf{«KR - коэффициент компенсации тока
3I0 по активному сопротивлению»}  и  \textbf{«KХ  -  коэффициент  компенсации тока  3I0 по реактивному сопротивлению»}.  

\item 
Коэффициенты компенсации рассчитываются по следующей формуле:
% ФОРМУЛЫ
\begin{align}
\vec K = \frac{1}{3} \cdot \frac{\vec Z_{0 \textsl{ уд}} - \vec Z_{1 \textsl{ уд}}}{\vec Z_{1 \textsl{ уд}}}\label{eq:groupK1},
\end{align}

\begin{eqexpl}
\item{$ \vec Z_{0 \textsl{ уд}} $}  комплексное удельное сопротивление нулевой последовательности;
\item{$ \vec Z_{1 \textsl{ уд}} $}  комплексное удельное сопротивление прямой последовательности.
\end{eqexpl}
\vspace{4pt}

\item 
Характеристики  срабатывания  ненаправленного  РС  для  междуфазных  и  фазных  контуров  задаются независимо для каждой ступени уставками:
\begin{itemize} 
\item \textbf{«Rср ДЗ-ФФ(ФЗ)-1(2..6)»} – активное сопротивление срабатывания РС; 
\item \textbf{«Хср ДЗ-ФФ(ФЗ)-1(2..6)»} – реактивное сопротивление срабатывания РС;
\item \textbf{«Ф1 ДЗ-ФФ(ФЗ)-1(2..6)»} – угол линии электропередач; 
\end{itemize}

\item 
Для  предотвращения  неселективного  срабатывания  РС-ФЗ-1  и  РС-ФФ-1  при  внешних  повреждениях
через  переходное  сопротивление  предусмотрена  возможность  отсечения  области  характеристики
срабатывания,  задаваемой  углом  $φ_4$.  Область  отсекается  отрезком,  исходящей  под  углом  $φ_4$,  из  точки  на
верхней части характеристики, соответствующей углу линии электропередачи  $φ_1$. Угол $φ_4$ задается уставкой
\textbf{«Ф4 ДЗ-ФФ(ФЗ)-1»}. 

\item 
Вид характеристики срабатывания ненаправленного РС представлен на рисунке \ref{pict_dz:char1}. 

\begin{figure}[H]
\centering
\includegraphics[width=0.4\textwidth,height=0.4\textheight,keepaspectratio]{\fbpath/051.0041 ДЗ/img/_char1.pdf}
\captionof{figure}{ Характеристика срабатывания ненаправленного РС}\label{pict_dz:char1}
\end{figure}

\item 
Параметры реле сопротивления приведены в таблице \ref{desc_dz:tbl1}

\small
\setlength{\extrarowheight}{0.05cm} % высота ячеек в таблице
\begin{longtable}
{|>{\centering\arraybackslash}m{15.2cm}|>{\centering\arraybackslash}m{1.7cm}|}
\caption{Параметры РС\hfill\vspace{-0.5\baselineskip}}\label{desc_dz:tbl1}\\ 
\hhline{|-|-|}
\rowcolor{gray!30}
Наименование параметра & Значение  \\ 
\hhline{|-|-|}
\endfirsthead

\caption*{\hspace{3pt}\emph{Продолжение таблицы \ref{desc_dz:tbl1}\hfill\vspace{-0.5\baselineskip}}} \\ % сделано по ГОСТ 2.105 п.6.8.7
\hhline{|-|-|}
\rowcolor{gray!30}
Наименование параметра & Значение  \\ 
\hhline{|-|-|}
\endhead
\hhline{|-|-|}

\endfoot
\endlastfoot

\raggedright Минимальное  междуфазное  напряжение,  при  котором  обеспечиваются  точностные характеристики РС, В, не менее & \centering\arraybackslash 0,5  \\  \hline
\raggedright Ток точной работы во всем диапазоне уставок при работе на угле $φ_1$, \% от номинального тока, не более  & \centering\arraybackslash 10  \\  \hline
\raggedright Собственное  время  срабатывания  РС  при  работе  на  угле  максимальной  чувствительности, трехкратном токе точной работы и скачкообразном уменьшении напряжения  на входе РС от напряжения,  соответствующего  сопротивлению  на  зажимах  РС  не  менее  $1,2\cdot(X_{\textsl{уст}} /sinφ_1 )$ до напряжения соответствующего $0,6\cdot(X_{\textsl{уст}} /sinφ_1 )$ мс, не более   & \centering\arraybackslash 25  \\  \hline
\raggedright Собственное  время  возврата  РС  при  работе  на  угле  максимальной  чувствительности, трехкратном  токе  точной  работы  и  скачкообразном  увеличении  напряжения  на  входе  РС  от напряжения соответствующего сопротивлению на зажимах РС $0,1\cdot(X_{\textsl{уст}} /sinφ_1 )$ до напряжения соответствующего $1,2\cdot(X_{\textsl{уст}} /sinφ_1 )$ (но не более 100 В для междуфазного канала и 57 В –- для фазного), мс, не более   & \centering\arraybackslash 50  \\  \hline
\raggedright Средняя основная погрешность всех РС по величине сопротивления срабатывания $R_{\textsl{уст}}$  и $X_{\textsl{уст}}$ при токе, равном $I_{\textsl{ном}}$ (или, в зависимости от уставки, меньшем токе, исходя из максимального напряжения на зажимах РС, равного 100 В), \% от уставки, не более  & \centering\arraybackslash $\pm$5  \\  \hline
\raggedright Основная абсолютная погрешность РС по углу наклона характеристики при токе, равном $I_{\textsl{ном}}$, эл. град., не более   & \centering\arraybackslash $\pm$5  \\  \hline
\raggedright Дополнительная абсолютная погрешность РС по углу наклона характеристики при изменении тока КЗ в диапазоне $(0,2 \dots 30) \cdot I_{\textsl{ном}}$, эл. град. от значения, измеренного при $I_{\textsl{ном}}$, не более   & \centering\arraybackslash $\pm$7  \\  \hline
\raggedright Дополнительная  погрешность  РС  по  сопротивлениям  $R$  и  $X$  при  изменении  температуры окружающей  среды  в  рабочем  диапазоне,  \%  от  среднего  значения,  измеренного  при нормальных климатических условиях, не более   & \centering\arraybackslash $\pm$3  \\  \hline
\raggedright Время  срабатывания  ДЗ  при  переходе  внешнего  КЗ  во  внутреннее  в  условиях  наличия насыщения ТТ мс, не более & \centering\arraybackslash 60  \\  \hline
\raggedright Коэффициент возврата, не более   & \centering\arraybackslash 1,1  \\  \hline

\end{longtable}
\setlength{\extrarowheight}{0.0cm} % возврат параметра высота ячеек в таблице
\normalsize

\item 
Выполнение функций направленного РС обеспечивается совместным использованием ненаправленного
РС и органов направленности (ОН). Орган направленности представляет собой два отрезка, исходящие из
начала координат и расположенные во втором и четвёртом квадрантах. Вид характеристики  (см. рис. \ref{pict_dz:on}) определяется
углами $φ_2$ и $φ_3$  наклона  этих  отрезков,  отсчитываемых  относительно  оси  $R$.  Углы  наклона  характеристики
задаются уставками \textbf{«Ф2 ОН-ФФ(ФЗ)»} и\textbf{«Ф3 ОН-ФФ(ФЗ)»}. Уставки органов направленности являются общими
для всех ступеней, задаются раздельно для контуров «фаза-фаза» и «фаза-земля». Углы наклона в ИЭУ
задаются  для  прямонаправленного  ОН в сторону АТ,  а  характеристика  обратнонаправленного  ОН  симметрична характеристике прямонаправленного ОН с поворотом на 180 градусов.

\begin{figure}[H]
\centering
\includegraphics[width=0.6\textwidth,height=0.6\textheight,keepaspectratio]{\fbpath/051.0041 ДЗ/img/_on.pdf}
\captionof{figure}{Характеристики срабатывания органов направления}\label{pict_dz:on}
\end{figure}

\item
Направленность для каждой ступени ДЗ задается с помощью программных переключателей \textbf{«Направл. ДЗ-ФФ(ФЗ)-1(2..6)»}, имеющих три положения:
\begin{enumerate}[label=\arabic*)]
\item \textit{Ненапр.} – ступень выполнена ненаправленной; 
\item \textit{В АТ} – ступень выполнена прямонаправленной в сторону АТ; 
\item \textit{В сеть} – ступень выполнена обратнонаправленной в сторону сети. 
\end{enumerate}

\item
При  близких  трехфазных  КЗ,  когда  все  напряжения  на  входе  РС  близки  к  нулю,  для  определения
направленности в течение времени не менее 40 мс используются напряжения предаварийного режима (работа
по «памяти»). При работе РС «по памяти» при трехфазных КЗ в месте установки защиты обеспечивается
длительность сигнала срабатывания на выходе РС не менее 40 мс в диапазоне токов от $0,2\cdot I_{\textsl{ном}}$ до $30\cdot I_{\textsl{ном}}$.
При КЗ «за спиной» с токами до $30\cdot I_{\textsl{ном}}$  обеспечивается отсутствие ложных срабатываний РС, в том числе при
снижении измеряемого напряжения до нуля. 

\item
Для отстройки от нагрузочного режима возможен вырез в характеристике срабатывания РС.
Предусмотрено два органа выреза нагрузочного режима для ступеней направленных в АТ и для ступеней
направленных в сеть, уставки органов выреза в АТ и в сеть для контуров «фаза-фаза» и «фаза-земля»
задаются раздельно.  Характеристика  срабатывания  приведена  на  рисунке \ref{pict_dz:r}.  Эта характеристика
срабатывания определяется уставками: 
\begin{itemize}
    \item \textbf{«Rнг-ФФ(ФЗ)в АТ»} - сопротивление выреза нагрузочного режима ступеней, направленных в сторону АТ;
    \item \textbf{«Фнг-ФФ(ФЗ) в АТ»} - угол выреза нагрузочного режима ступеней, направленных в сторону АТ;
	\item \textbf{«Rнг-ФФ(ФЗ) в сеть»} - сопротивление выреза нагрузочного режима ступеней, направленных в сторону сети;
	\item \textbf{«Фнг-ФФ(ФЗ) в сеть»} - угол выреза нагрузочного режима ступеней, направленных в сторону сети.
\end{itemize}

\item
Исключаемая область в характеристике срабатывания симметрична относительно осей координат. Ввод
в работу отстройки от нагрузки в сторону АТ и в сторону сети осуществляется с помощью программных
переключателей \textbf{«Zнг в АТ»} и \textbf{«Zнг в сеть»}  соответственно. 

\begin{figure}[H]
\centering
\includegraphics[width=0.5\textwidth,height=0.5\textheight,keepaspectratio]{\fbpath/051.0041 ДЗ/img/_r.pdf}
\captionof{figure}{Характеристика срабатывания измерительного органа отстройки от нагрузки}\label{pict_dz:r}
\end{figure}

\item
Суммарная характеристика направленного РС, используемого в логике работы ступеней ДЗ приведена на рисунке \ref{pict_dz:rs}.  

\begin{figure}[H]
\centering
\includegraphics[width=0.75\textwidth,height=0.75\textheight,keepaspectratio]{\fbpath/051.0041 ДЗ/img/_rs.pdf}
\captionof{figure}{Характеристики срабатывания направленного РС}\label{pict_dz:rs}
\end{figure}

\item
Для  обнаружения  однофазных  КЗ  и  пуска  ступеней  ДЗ-ФЗ,  контролирующих  контуры  «фаза-земля»,
предусмотрен  орган  определения  вида  повреждения  (ООВП),  основанный  на  совместном  применении
измерительного органа тока нулевой последовательности (ИО 3I0 ООВП) с торможением от одного из фазных
токов (средний из  $I_A , I_B , I_C$) и измерительного органа напряжения нулевой последовательности (ИО 3U0 ООВП).
Торможение осуществляется от тока той фазы, значение в которой является средним между максимальным и
минимальным значениями токов в остальных двух фазах. В таком случае при междуфазном КЗ на землю
торможение максимальном, а при однофазном КЗ –- минимально. 

\item
Характеристика  срабатывания  ИО  3I0  ООВП  с  торможением  представлена  на  рисунке  \ref{pict_dz:torm}.  

\begin{figure}[H]
\centering
\includegraphics[width=0.65\textwidth,height=0.65\textheight,keepaspectratio]{\fbpath/051.0041 ДЗ/img/_torm.pdf}
\captionof{figure}{Характеристика срабатывания ИО 3I0 ООВП}\label{pict_dz:torm}
\end{figure}

\item
Порог срабатывания горизонтального участка задается уставкой \textbf{«3I0ср ООВП»}. Наклонный участок
характеристики проходит через точку, равной $1,25 \cdot I_{\textsl{ном}}$ , под углом к горизонтальной оси, соответствующим
коэффициенту торможения $k_T$. Коэффициент торможения задается уставкой \textbf{«Кт ООВП»}. При превышении
тормозного тока уставки \textbf{«Iср блок. торм.»} ООВП блокируется.

\begin{figure}[H]
\centering
\includegraphics[width=0.65\textwidth,height=0.65\textheight,keepaspectratio]{\fbpath/051.0041 ДЗ/img/OOVP.pdf}
\captionof{figure}{Функциональная схема ООВП}\label{pict_dz:OOVP}
\end{figure}

\item
Ввод  в  работу  ступеней  ДЗ,  работающих  по  контуру  «фаза-земля»  может  осуществляться  с
дополнительным  контролем  напряжения  нулевой  последовательности  за  счет  ввода  программного
ключа \textbf{«Контр.  3U0 ООВП»}.  Напряжение  срабатывания  ИО  3U0  ООВП  задается  уставкой  \textbf{«3U0ср
ООВП»}. ИО 3U0 ООВП может работать по напряжению $3U_{\textsl{0 звезда}}$, рассчитанному на основе фазных 
напряжений, или по напряжению $3U_{\textsl{0 треуг.}}$   --  выбирается  программным  ключом  \textbf{«Выбор  3U0»},
имеющим  два  соответствующих положения: \textsl{«Расчетн.»} и \textsl{«Измерен.»}.

\item
Переменная $3U_{\textsl{0 треуг.}}$ зависит от положения программного переключателя \textbf{«Схема подключения цепей
разомкнутого треугольника»}.  Данный  переключатель  позволяет  учесть  различные  схемы  соединения
дополнительной  обмотки  ТН.  Если  переключатель  выставлен  в  положение  «$U_{\textsl{НК}}$»,  то  переменной  $3U_{\textsl{0 треуг.}}$
присваивается  значение  напряжения,  измеренного  по  аналоговому  входу  «$U_{\textsl{ИФ}}(U_{\textsl{НК}})$».  Если  программный
переключатель {\textbf{«Схема подключения цепей треугольника»}} выставлен в положение «$U_{\textsl{НИ}}/U_{\textsl{ИК}}$», «$U_{\textsl{НФ}}/U_{\textsl{ФК}}$»
или  «$U_{\textsl{НИ}}/U_{\textsl{ИФ}}/U_{\textsl{ФК}}$», то переменная $3U_{\textsl{0 треуг.}}$ рассчитывается в соответствии с формулами, приведенными в
приложении \ref{app:signals}. Программный переключатель \textbf{«Выбор 3U0»} аналогично влияет на работу ОНМНП. 

\item
В составе логики дистанционной защиты предусмотрено шесть ступеней от междуфазных замыканий с
контролем трех контуров «фаза-фаза». 

\vspace{2mm}
\begin{figure}[H]
\centering
\includegraphics[width=0.65\textwidth,height=0.65\textheight,keepaspectratio]{\fbpath/051.0041 ДЗ/img/DZ_FF_1.pdf}
\captionof{figure}{Функциональная схема ДЗ-ФФ-1}\label{pict_dz:DZ_FF_1}
\end{figure}

\item
Пороги срабатывания РС ступеней ДЗ-ФФ задаются уставками \textbf{«Хср-ФФ-1(2…6)»}, \textbf{«Rср-ФФ-1(2…6)»} и
\textbf{«Ф1-ФФ-1(2…6)»}, а выдержки времени срабатывания ступеней регулируются уставками \textbf{«Тср ДЗ-ФФ-1(2…6)»}.
Ввод в работу ступеней ДЗ-ФФ осуществляется с помощью программных переключателей \textbf{«ДЗ-ФФ-1(2…6)»}.
Направленности для каждой ступени задается с помощью программных переключателей \textbf{«Направл ДЗ-ФФ-1(2…6)»}.

\item
Предусмотрена возможность ввода подхвата направленного пуска ступеней ДЗ-ФФ от ненаправленного
пуска той же ступени, задается программным переключателем \textbf{«Подхват РС-ФФ-1(2…6)»}. Тем самым
гарантируется правильное действие защиты при близких замыканиях. 

\item
Оперативный ввод/вывод ДЗ-ФФ производится с помощью БУ (блоков управления функциями) \textit{«Вывод ДЗ»}.

\item
Пуски степеней ДЗ-ФФ направленные в сеть, действуют в логику предварительного деления своей
стороны АТ. С помощью выдержек времени в логике деления \textbf{«Тср ДЗ-ФФ-1(2…6) на СВ»}, \textbf{«ДЗ-ФФ-1(2…6) на
В сеч»} и \textbf{«Тср ДЗ-ФФ-1(2…6) на В»}, первым действием формируются сигналы на разделение секции
(отключение СВ) или деление по сечению (отключение В сечения), вторым действием формируются сигналы
на отключение выключателей своей стороны с пуском УРОВ и запретом АПВ. Пуски ступеней ДЗ-ФФ
направленные в АТ, могут действовать в логику предварительного деления резервных защит смежной стороны
(смежный комплект РЗАТ), через выдержки времени в логике деления \textbf{«Тср ДЗ-ФФ-1(2…6) на пуск ЛД см ст»}.
Действие срабатывания ступеней ДЗ-ФФ предусматривает отключение АТ со всех сторон с запретом АПВ и
пуском УРОВ.

\item
В составе логики дистанционной защиты предусмотрено две ступеней от однофазных замыканий с
контролем трех контуров «фаза-земля».

\vspace{2mm}
\begin{figure}[H]
\centering
\includegraphics[width=0.65\textwidth,height=0.65\textheight,keepaspectratio]{\fbpath/051.0041 ДЗ/img/DZ_FZ_1.pdf}
\captionof{figure}{Функциональная схема ДЗ-ФЗ-1}\label{pict_dz:DZ_FZ_1}
\end{figure}

\item
Пуск ступеней ДЗ-ФЗ разрешается только при срабатывании органа определения вида повреждения
\textit{«Ввод ДЗ-ФЗ»}. Пороги срабатывания РС ступеней ДЗ-ФЗ задаются уставками \textbf{«Хср-ФЗ-1(2)»}, \textbf{«Rср-ФЗ-1(2)»} и
\textbf{«Ф1-ФЗ-1(2)»}, а выдержки времени срабатывания ступеней регулируются уставками \textbf{«Тср ДЗ-ФЗ-1(2)»}. Ввод в
работу ступеней ДЗ-ФЗ осуществляется с помощью программных переключателей \textbf{«ДЗ-ФЗ-1(2)»}.
Направленности для каждой ступени задается с помощью программных переключателей \textbf{«Направл ДЗ-ФФ-1(2)»}.

\item
Предусмотрена возможность ввода подхвата направленного пуска ступеней ДЗ-ФЗ от ненаправленного
пуска той же ступени, задается программными переключателями \textbf{«Подхват РС-ФЗ-1(2)»}. Тем самым
гарантируется правильное действие защиты при близких замыканиях. 

\item
Оперативный ввод/вывод ДЗ-ФЗ производится с помощью БУ (блока управления функциями) \textit{«Вывод ДЗ»}.

\item
Пуски степеней ДЗ-ФЗ направленные в сеть, действуют в логику предварительного деления своей
стороны АТ. С помощью выдержек времени в логике деления \textbf{«Тср ДЗ-ФЗ-1(2) на СВ»}, \textbf{«ДЗ-ФФ-1(2) на В сеч»}
и \textbf{«Тср ДЗ-ФФ-1(2) на В»}, первым действием формируются сигналы на разделение секции (отключение СВ)
или деление по сечению (отключение В сечения), вторым действием формируются сигналы на отключение
выключателей своей стороны с пуском УРОВ и запретом АПВ. Действие срабатывания ступеней ДЗ-ФЗ
предусматривает отключение АТ со всех сторон с запретом АПВ и пуском УРОВ.

\item
Для исключения ложного действия при возникновении неисправности в цепях напряжения ступени ДЗ
могут быть заблокированы от БНН. Блокировка от БНН вводится отдельно для каждой ступени ДЗ с помощью
программных переключателей \textbf{«БНН ДЗ-ФФ-1(2…6)»} и \textbf{«БНН ДЗ-ФЗ-1(2)»}.

\item
Дистанционная защита контролируется блокировкой при качаниях, предотвращающей ложное
срабатывание ДЗ при качаниях и асинхронном ходе. В устройстве реализовано два способа блокировки:
\begin{itemize}
\item по приращению токов прямой и обратной последовательности;
\item по скорости изменения величины сопротивления dZ/dt.
\end{itemize}

\item
Контроль ступеней ДЗ от блокировки при качаниях вводится программными переключателями
\textbf{«БК ДЗ-ФФ(ФЗ)-1(2…6)»}, имеющими четыре положения:
\begin{itemize}
\item \textit{«Выведено»} -- ступень работает без контроля от БК;
\item \textit{«БК-I-б»} -- ступень пускается от сигнала ввода быстродействующих ступеней;
\item \textit{«БК-I-м»} -- ступень пускается от сигнала ввода медленнодействующих ступеней;
\item \textit{«БК-Z»} -- ступень пускается от сигнала БК по скорости изменения величины сопротивления.
\end{itemize}

\item
Если за время пуска от БК происходит срабатывание ненаправленного РС ступени ДЗ, то разращение
от БК подхватывается и удерживается до пропадания условий срабатывания ненаправленного РС
соответствующей ступени.

\item
В составе логики дистанционной защиты предусмотрено четыре режима оперативного ускорения ДЗ.





































\end{enumerate}

%%%%%%%%%%%%%%%%%%%%%%%%%%%%%%%%%%%%%%%%%%%%%%%%%%%%%%%%%%%%%%%%%%%%%% ПЕРВАЯ СТУПЕНЬ %%%%%%%%%%%%%%%%%%%%%%%%%%%%%%%%%%%%%%%%%%%%%%%%%%
\item\label{sec:sec2} Первая ступень ДЗ от междуфазных КЗ (ДЗ-ФФ-1)
\begin{enumerate}[label=\arabic{section}.\arabic{subsection}.\arabic{enumi}.\arabic*, labelsep=4pt, leftmargin=0em, itemindent=65pt, parsep=0pt]

\item
Первая  ступень  ДЗ  от  междуфазных  КЗ  содержит  две  «подступени»:  быстродействующая ДЗ-ФФ-1(б)  и  медленнодействующая  ДЗ-ФФ-1(м),  которая  вводится программным  ключом  «ДЗ-ФФ-1(м)». Всю ступень можно ввести/вывести программным ключом \textbf{«ДЗ-ФФ-1»}. Контроль от блокировки при качаниях (БК)  для  каждой  подступени  задается  отдельными  программными  переключателями  \textbf{«БК  ДЗ-ФФ-1(б)»}  и  \textbf{«БК  ДЗ-ФФ-1(м)»}.  Действие  подступеней  осуществляется  с  регулируемыми  независимыми  выдержками времени \textbf{«Тср ДЗ-ФФ-1(б)»} и \textbf{«Тср ДЗ-ФФ-1(м)»}.


\item
Если междуфазное КЗ произошло в зоне действия ДЗ-ФФ-1 и при этом в течение времени ввода БК сработало направленное РС-ФФ-2 с охватом начала координат, то разрешающий сигнал от БК подхватывается и удерживается независимо от истечения времени ввода. Возврат разрешающего сигнала происходит только после возврата направленного РС-ФФ-2 с охватом начала координат. 

\item
Функциональная схема логики первой ступени  ДЗ  от  междуфазных  КЗ приведена на рисунке \ref{pict_dz:ff1}.


\item
Уставки первой ступени ДЗ от междуфазных КЗ (ДЗ-ФФ-1) представлены в таблице \ref{desc_dz:tbl3}.

\small
\begin{longtable}{|>{\centering\arraybackslash}m{5.3cm}|>{\centering\arraybackslash}m{3.3cm}|>{\centering\arraybackslash}m{4.2cm}|>{\centering\arraybackslash}m{1.8cm}|>{\centering\arraybackslash}m{1cm}|}
% Первая страница: полный заголовок
\caption{Параметры для настройки ступени ДЗ-ФФ-1\hfill\vspace{-0.5\baselineskip}}\label{desc_dz:tbl3}\\
\hline
\rowcolor{gray!30}
Параметр & Условное обозначение на схеме & Значение/ Диапазон & Единица измерения & Шаг \\
\hline
\endfirsthead

% Промежуточные страницы: «Продолжение таблицы»
\caption*{\hspace{3pt}\emph{Продолжение таблицы \ref{desc_dz:tbl3}}\hfill\ }\vspace{-0.5\baselineskip} \\
\hline
\rowcolor{gray!30}
Параметр & Условное обозначение на схеме & Значение/ Диапазон & Единица измерения & Шаг \\
\hline
\endhead

%===m>ДЗ|ДЗ-ФФ-1
\raggedright {'Description': ''} & \centering ДЗ-ФФ-1 & \centering Выведено \\ Введено & \centering --- & \centering \arraybackslash --- \\
\hline
\raggedright {'Description': ''} & \centering RelesoprotivleniyadlyaDZFF11Rsr & \centering 0,000 ... 100,000 & \centering  & \centering \arraybackslash 0,001 \\
\hline
\raggedright {'Description': ''} & \centering RelesoprotivleniyadlyaDZFF11Xsr & \centering 0,000 ... 100,000 & \centering  & \centering \arraybackslash 0,001 \\
\hline
\raggedright {'Description': ''} & \centering RelesoprotivleniyadlyaDZFF11F1 & \centering 0,000 ... 90,000 & \centering  & \centering \arraybackslash 0,001 \\
\hline
\raggedright {'Description': ''} & \centering RelesoprotivleniyadlyaDZFF11F4 & \centering 0,000 ... 20,000 & \centering  & \centering \arraybackslash 0,001 \\
\hline
\raggedright {'Description': ''} & \centering RelesoprotivleniyadlyaDZFF11F2 & \centering 0,000 ... 359,000 & \centering  & \centering \arraybackslash 0,001 \\
\hline
\raggedright {'Description': ''} & \centering Тср ДЗ-ФФ (б) & \centering 0 ... 100000 & \centering  & \centering \arraybackslash 1 \\
\hline
\raggedright {'Description': ''} & \centering RelesoprotivleniyadlyaDZFF11F3 & \centering 0,000 ... 359,000 & \centering  & \centering \arraybackslash 0,001 \\
\hline
\raggedright {'Description': ''} & \centering RelesoprotivleniyadlyaDZFF11NapravlennostDZFF & \centering Прямо \\ Обратно & \centering --- & \centering \arraybackslash --- \\
\hline
\raggedright {'Description': ''} & \centering Подхват РС-ФФ & \centering Введено \\ Выведено & \centering --- & \centering \arraybackslash --- \\
\hline
\raggedright {'Description': ''} & \centering RelesoprotivleniyadlyaDZFF11Rng & \centering 0,000 ... 100,000 & \centering  & \centering \arraybackslash 0,001 \\
\hline
\raggedright {'Description': ''} & \centering RelesoprotivleniyadlyaDZFF11Fng & \centering 0,000 ... 45,000 & \centering  & \centering \arraybackslash 0,001 \\
\hline
\raggedright {'Description': ''} & \centering Тср ДЗ-ФФ (м) & \centering 0 ... 100000 & \centering  & \centering \arraybackslash 1 \\
\hline
\raggedright {'Description': ''} & \centering ДЗ-ФФ(м) & \centering Введено \\ Выведено & \centering --- & \centering \arraybackslash --- \\
\hline
\raggedright {'Description': ''} & \centering PereklyuchatelVkh4boolconst1S & \centering Откл. \\ БК-I-б \\ БК-I-м \\ БК-Z & \centering --- & \centering \arraybackslash --- \\
\hline
\raggedright {'Description': ''} & \centering PereklyuchatelVkh4boolconst2S & \centering Откл. \\ БК-I-б \\ БК-I-м \\ БК-Z & \centering --- & \centering \arraybackslash --- \\
\hline
\raggedright {'Description': ''} & \centering БСТО ДЗ-ФФ & \centering Введено \\ Выведено & \centering --- & \centering \arraybackslash --- \\
\hline
\raggedright {'Description': ''} & \centering БНН ДЗ-ФФ & \centering Введено \\ Выведено & \centering --- & \centering \arraybackslash --- \\
\hline
%===m
\end{longtable}
\normalsize

\end{enumerate}
%%%%%%%%%%%%%%%%%%%%%%%%%%%%%%%%%%%%%%%%%%%%%%%%%%%%%%%%%%%%%%%%%%%%%% КОНЕЦ ПЕРВАЯ СТУПЕНЬ %%%%%%%%%%%%%%%%%%%%%%%%%%%%%%%%%%%%%%%%%%%%%%%%%%


%%%%%%%%%%%%%%%%%%%%%%%%%%%%%%%%%%%%%%%%%%%%%%%%%%%%%%%%%%%%%%%%%%%%%% ВТОРАЯ СТУПЕНЬ %%%%%%%%%%%%%%%%%%%%%%%%%%%%%%%%%%%%%%%%%%%%%%%%%%
\item\label{sec:sec2} Вторая ступень ДЗ от междуфазных КЗ (ДЗ-ФФ-2)
\begin{enumerate}[label=\arabic{section}.\arabic{subsection}.\arabic{enumi}.\arabic*, labelsep=4pt, leftmargin=0em, itemindent=65pt, parsep=0pt]

\item
Вторая ступень ДЗ от междуфазных КЗ также выполнена с двумя выдержками времени и с раздельным контролем от БК. Вся ступень вводится программным ключом \textbf{«ДЗ-ФФ-2»}. Быстродействующую подступень можно  ввести  в  работу  программным  ключом  \textbf{«ДЗ-ФФ-2(б)»}.  Действие  подступеней  осуществляется  с регулируемыми независимыми выдержками времени \textbf{«Тср ДЗ-ФФ-2(б)»} и \textbf{«Тср ДЗ-ФФ-2(м)»}. Контроль от БК для каждой подступени задается отдельными программными переключателями \textbf{«БК ДЗ-ФФ-2(б)»} и \textbf{«БК ДЗ-ФФ-2(м)»}.

\item
Разрешающий сигнал от БК для ДЗ-ФФ-2 подхватывается от направленного РС-ФФ-2 с охватом начала координат  и  удерживается  независимо  от  истечения  времени  ввода  БК.  Возврат  разрешающего  сигнала происходит только после возврата направленного РС-ФФ-2 с охватом начала координат.

\item
Функциональная схема логики второй ступени ДЗ от междуфазных КЗ приведена на рисунке \ref{pict_dz:ff2}.

\vspace{2mm}
\begin{figure}[H]
\centering
\includegraphics[width=1\textwidth,height=1\textheight,keepaspectratio]{\fbpath/051.0041 ДЗ/img/DZ_FF_2.pdf}
\captionof{figure}{Функциональная схема логики ДЗ-ФФ-2}\label{pict_dz:ff2}
\end{figure}

\item
Уставки второй ступени ДЗ от междуфазных КЗ (ДЗ-ФФ-2) представлены в таблице \ref{desc_dz:tbl4}.

\small
\begin{longtable}{|>{\centering\arraybackslash}m{5.3cm}|>{\centering\arraybackslash}m{3.3cm}|>{\centering\arraybackslash}m{4.2cm}|>{\centering\arraybackslash}m{1.8cm}|>{\centering\arraybackslash}m{1cm}|}
% Первая страница: полный заголовок
\caption{Параметры для настройки ступени ДЗ-ФФ-2\hfill\vspace{-0.5\baselineskip}}\label{desc_dz:tbl4}\\
\hline
\rowcolor{gray!30}
Параметр & Условное обозначение на схеме & Значение/ Диапазон & Единица измерения & Шаг \\
\hline
\endfirsthead

% Промежуточные страницы: «Продолжение таблицы»
\caption*{\hspace{3pt}\emph{Продолжение таблицы \ref{desc_dz:tbl4}}\hfill\ }\vspace{-0.5\baselineskip} \\
\hline
\rowcolor{gray!30}
Параметр & Условное обозначение на схеме & Значение/ Диапазон & Единица измерения & Шаг \\
\hline
\endhead

%===m>ДЗ|ДЗ-ФФ-2
\raggedright {'Description': ''} & \centering Rср & \centering 0,000 ... 100,000 & \centering  & \centering \arraybackslash 0,001 \\
\hline
\raggedright {'Description': ''} & \centering Xср & \centering 0,000 ... 100,000 & \centering  & \centering \arraybackslash 0,001 \\
\hline
\raggedright {'Description': ''} & \centering Ф1 & \centering 0,000 ... 359,000 & \centering  & \centering \arraybackslash 0,001 \\
\hline
\raggedright {'Description': ''} & \centering Ф2 & \centering 0,000 ... 359,000 & \centering  & \centering \arraybackslash 0,001 \\
\hline
\raggedright {'Description': ''} & \centering Ф3 & \centering 0,000 ... 359,000 & \centering  & \centering \arraybackslash 0,001 \\
\hline
\raggedright {'Description': ''} & \centering Направленность ДЗ-ФФ & \centering Прямо \\ Обратно & \centering --- & \centering \arraybackslash --- \\
\hline
\raggedright {'Description': ''} & \centering Rнг & \centering 0,000 ... 100,000 & \centering  & \centering \arraybackslash 0,001 \\
\hline
\raggedright {'Description': ''} & \centering Фнг & \centering 0,000 ... 359,000 & \centering  & \centering \arraybackslash 0,001 \\
\hline
\raggedright {'Description': ''} & \centering ДЗ-ФФ-2 & \centering Выведено \\ Введено & \centering --- & \centering \arraybackslash --- \\
\hline
\raggedright {'Description': ''} & \centering ДЗ-ФФ-2(б) & \centering Введено \\ Выведено & \centering --- & \centering \arraybackslash --- \\
\hline
\raggedright {'Description': ''} & \centering NezavisimoeRVnasrabatyvanie1Tvremya & \centering 0 ... 100000 & \centering  & \centering \arraybackslash 1 \\
\hline
\raggedright {'Description': ''} & \centering БК-ДЗ-ФФ-2 (б) & \centering Откл. \\ БК-I-б \\ БК-I-м \\ БК-Z & \centering --- & \centering \arraybackslash --- \\
\hline
\raggedright {'Description': ''} & \centering NezavisimoeRVnasrabatyvanie2Tvremya & \centering 0 ... 100000 & \centering  & \centering \arraybackslash 1 \\
\hline
\raggedright {'Description': ''} & \centering БК-ДЗ-ФФ-2 (м) & \centering Откл. \\ БК-I-б \\ БК-I-м \\ БК-Z & \centering --- & \centering \arraybackslash --- \\
\hline
\raggedright {'Description': ''} & \centering БНН ДЗ-ФФ-2 & \centering Введено \\ Выведено & \centering --- & \centering \arraybackslash --- \\
\hline
%===m
\end{longtable}
\normalsize

\end{enumerate}

\newpage
%%%%%%%%%%%%%%%%%%%%%%%%%%%%%%%%%%%%%%%%%%%%%%%%%%%%%%%%%%%%%%%%%%%%%% КОНЕЦ ВТОРАЯ СТУПЕНЬ %%%%%%%%%%%%%%%%%%%%%%%%%%%%%%%%%%%%%%%%%%%%%%%%%%


%%%%%%%%%%%%%%%%%%%%%%%%%%%%%%%%%%%%%%%%%%%%%%%%%%%%%%%%%%%%%%%%%%%%%% 3,4,5,6 СТУПЕНИ %%%%%%%%%%%%%%%%%%%%%%%%%%%%%%%%%%%%%%%%%%%%%%%%%%
\item\label{sec:sec21} Остальные ступени ДЗ от междуфазных КЗ (ДЗ-ФФ-3, ДЗ-ФФ-4, ДЗ-ФФ-5, ДЗ-ФФ-6)
\begin{enumerate}[label=\arabic{section}.\arabic{subsection}.\arabic{enumi}.\arabic*, labelsep=4pt, leftmargin=0em, itemindent=65pt, parsep=0pt]

\item
Логика  работы  остальных  ступеней  ДЗ  от  междуфазных  КЗ  выполнена  одинаковой.  Ввод  в  работу ступеней осуществляется программными ключами \textbf{«ДЗ-ФФ-3(4,5,6)»}. Действие ступеней осуществляется с регулируемыми  независимыми  выдержками  времени  \textbf{«Тср  ДЗ-ФФ-3(4,5,6)»}.  Контроль  от  БК  задается отдельными программными переключателями \textbf{«БК ДЗ-ФФ-3(4,5,6)»}.  

\item
Разрешающий  сигнал  от  БК  для  ступеней  подхватывается  от  срабатывания  РС  соответствующей ступени  и  удерживается  независимо  от  истечения  времени  ввода  БК.  Возврат  разрешающего  сигнала происходит только после возврата РС соответствующей ступени. На рисунке \ref{pict_dz:ff3} приведена функциональная схема логики ДЗ-ФФ-3, выполнение функциональных схем логики ступеней ДЗ ДЗ-ФФ-4, ДЗ-ФФ-5, ДЗ-ФФ-6 аналогичное.

\vspace{2mm}
\begin{figure}[H]
\centering
\includegraphics[width=1\textwidth,height=1\textheight,keepaspectratio]{\fbpath/051.0041 ДЗ/img/DZ_FF_3.pdf}
\captionof{figure}{Функциональная схема логики ДЗ-ФФ-3}\label{pict_dz:ff3}
\end{figure}

\item
Уставки ступеней ДЗ от междуфазных КЗ (ДЗ-ФФ-3, ДЗ-ФФ-4, ДЗ-ФФ-5, ДЗ-ФФ-6) представлены в таблице \ref{desc_dz:tbl5}.

\small
\begin{longtable}{|>{\centering\arraybackslash}m{5.3cm}|>{\centering\arraybackslash}m{3.3cm}|>{\centering\arraybackslash}m{4.2cm}|>{\centering\arraybackslash}m{1.8cm}|>{\centering\arraybackslash}m{1cm}|}
% Первая страница: полный заголовок
\caption{Параметры для настройки ступени ДЗ-ФФ-3(4,5,6)\hfill\vspace{-0.5\baselineskip}}\label{desc_dz:tbl5}\\
\hline
\rowcolor{gray!30}
Параметр & Условное обозначение на схеме & Значение/ Диапазон & Единица измерения & Шаг \\
\hline
\endfirsthead

% Промежуточные страницы: «Продолжение таблицы»
\caption*{\hspace{3pt}\emph{Продолжение таблицы \ref{desc_dz:tbl5}}\hfill\ }\vspace{-0.5\baselineskip} \\
\hline
\rowcolor{gray!30}
Параметр & Условное обозначение на схеме & Значение/ Диапазон & Единица измерения & Шаг \\
\hline
\endhead

%===m>ДЗ|ДЗ-ФФ-3
\raggedright {'Description': ''} & \centering Rср & \centering 0,000 ... 100,000 & \centering  & \centering \arraybackslash 0,001 \\
\hline
\raggedright {'Description': ''} & \centering Xср & \centering 0,000 ... 100,000 & \centering  & \centering \arraybackslash 0,001 \\
\hline
\raggedright {'Description': ''} & \centering Ф1 & \centering 0,000 ... 359,000 & \centering  & \centering \arraybackslash 0,001 \\
\hline
\raggedright {'Description': ''} & \centering Ф2 & \centering 0,000 ... 359,000 & \centering  & \centering \arraybackslash 0,001 \\
\hline
\raggedright {'Description': ''} & \centering Ф3 & \centering 0,000 ... 359,000 & \centering  & \centering \arraybackslash 0,001 \\
\hline
\raggedright {'Description': ''} & \centering Направленность ДЗ-ФФ & \centering Прямо \\ Обратно & \centering --- & \centering \arraybackslash --- \\
\hline
\raggedright {'Description': ''} & \centering Rнг & \centering 0,000 ... 100,000 & \centering  & \centering \arraybackslash 0,001 \\
\hline
\raggedright {'Description': ''} & \centering Фнг & \centering 0,000 ... 359,000 & \centering  & \centering \arraybackslash 0,001 \\
\hline
\raggedright {'Description': ''} & \centering ДЗ-ФФ & \centering Выведено \\ Введено & \centering --- & \centering \arraybackslash --- \\
\hline
\raggedright {'Description': ''} & \centering Тср ДЗ-ФФ & \centering 0 ... 100000 & \centering  & \centering \arraybackslash 1 \\
\hline
\raggedright {'Description': ''} & \centering БК ДЗ-ФФ & \centering Откл. \\ БК-I-б \\ БК-I-м \\ БК-Z & \centering --- & \centering \arraybackslash --- \\
\hline
\raggedright {'Description': ''} & \centering БНН ДЗ-ФФ & \centering Введено \\ Выведено & \centering --- & \centering \arraybackslash --- \\
\hline
%===m
\end{longtable}
\normalsize

\end{enumerate}
%%%%%%%%%%%%%%%%%%%%%%%%%%%%%%%%%%%%%%%%%%%%%%%%%%%%%%%%%%%%%%%%%%%%%% КОНЕЦ 3,4,5,6 СТУПЕНИ %%%%%%%%%%%%%%%%%%%%%%%%%%%%%%%%%%%%%%%%%%%%%%%%%%


%%%%%%%%%%%%%%%%%%%%%%%%%%%%%%%%%%%%%%%%%%%%%%%%%%%%%%%%%%%%%%%%%%%%%% ПЕРВАЯ СТУПЕНЬ ФАЗА-ЗЕМЛЯ %%%%%%%%%%%%%%%%%%%%%%%%%%%%%%%%%%%%%%%%%%%%%%%%%%
\item\label{sec:sec1fz} Первая ступень ДЗ от однофазных КЗ (ДЗ-ФЗ-1)
\begin{enumerate}[label=\arabic{section}.\arabic{subsection}.\arabic{enumi}.\arabic*, labelsep=4pt, leftmargin=0em, itemindent=65pt, parsep=0pt]

\item
Первая ступень ДЗ от однофазных КЗ по аналогии с ДЗ-ФФ-1 выполнена с двумя выдержками времени и с раздельным контролем от БК. Вся ступень вводится программным ключом \textbf{«ДЗ-ФЗ-1»}. Медленнодействующую подступень можно ввести в работу программным ключом \textbf{«ДЗ-ФЗ-1(м)»}. Действие подступеней осуществляется с регулируемыми независимыми выдержками времени \textbf{«Тср ДЗ-ФЗ-1(б)»} и \textbf{«Тср ДЗ-ФЗ-1(м)»}. Контроль от БК для каждой подступени задается отдельными программными переключателями \textbf{«БК ДЗ-ФЗ-1(б)»} и \textbf{«БК ДЗ-ФЗ-1(м)»}.

\item
Предусмотрена возможность подхвата сработанного состояния РС-ФЗ-1 от направленного РС-ФЗ-1 с охватом начала координат (вводится программным ключом \textbf{«Подхват РС-ФЗ-1»}). Тем самым гарантируется правильное действие защиты при близких КЗ. Возврат схемы подхвата в исходное состояние происходит только после возврата направленного РС-ФЗ-1 с охватом начала координат.

\item
Разрешающий сигнал от БК для ДЗ-ФЗ-1 подхватывается от направленного РС-ФЗ-1 с охватом начала координат и удерживается независимо от истечения времени ввода БК. Возврат разрешающего сигнала происходит только после возврата направленного РС-ФФ-1 с охватом начала координат.

\item
Функциональная схема логики первой ступени ДЗ от однофазных КЗ приведена на рисунке \ref{pict_dz:DZ_FZ_1}.



\item
Уставки первой ступени ДЗ от однофазных КЗ (ДЗ-ФЗ-1) представлены в таблице \ref{desc_dz:tbl6}.

\small
\begin{longtable}{|>{\centering\arraybackslash}m{5.3cm}|>{\centering\arraybackslash}m{3.3cm}|>{\centering\arraybackslash}m{4.2cm}|>{\centering\arraybackslash}m{1.8cm}|>{\centering\arraybackslash}m{1cm}|}
% Первая страница: полный заголовок
\caption{Параметры для настройки ступени ДЗ-ФЗ-1\hfill\vspace{-0.5\baselineskip}}\label{desc_dz:tbl6}\\
\hline
\rowcolor{gray!30}
Параметр & Условное обозначение на схеме & Значение/ Диапазон & Единица измерения & Шаг \\
\hline
\endfirsthead

% Промежуточные страницы: «Продолжение таблицы»
\caption*{\hspace{3pt}\emph{Продолжение таблицы \ref{desc_dz:tbl6}}\hfill\ }\vspace{-0.5\baselineskip} \\
\hline
\rowcolor{gray!30}
Параметр & Условное обозначение на схеме & Значение/ Диапазон & Единица измерения & Шаг \\
\hline
\endhead

%===m>ДЗ|ДЗ-ФЗ-1
\raggedright {'Description': ''} & \centering RelesoprotivleniyadlyaDZFF11Rsr & \centering 0,000 ... 100,000 & \centering  & \centering \arraybackslash 0,001 \\
\hline
\raggedright {'Description': ''} & \centering RelesoprotivleniyadlyaDZFF11Xsr & \centering 0,000 ... 100,000 & \centering  & \centering \arraybackslash 0,001 \\
\hline
\raggedright {'Description': ''} & \centering RelesoprotivleniyadlyaDZFF11F1 & \centering 0,000 ... 90,000 & \centering  & \centering \arraybackslash 0,001 \\
\hline
\raggedright {'Description': ''} & \centering RelesoprotivleniyadlyaDZFF11F4 & \centering 0,000 ... 20,000 & \centering  & \centering \arraybackslash 0,001 \\
\hline
\raggedright {'Description': ''} & \centering RelesoprotivleniyadlyaDZFF11F2 & \centering 0,000 ... 359,000 & \centering  & \centering \arraybackslash 0,001 \\
\hline
\raggedright {'Description': ''} & \centering RelesoprotivleniyadlyaDZFF11F3 & \centering 0,000 ... 359,000 & \centering  & \centering \arraybackslash 0,001 \\
\hline
\raggedright {'Description': ''} & \centering RelesoprotivleniyadlyaDZFF11NapravlennostDZFF & \centering Прямо \\ Обратно & \centering --- & \centering \arraybackslash --- \\
\hline
\raggedright {'Description': ''} & \centering RelesoprotivleniyadlyaDZFF11Rng & \centering 0,000 ... 100,000 & \centering  & \centering \arraybackslash 0,001 \\
\hline
\raggedright {'Description': ''} & \centering RelesoprotivleniyadlyaDZFF11Fng & \centering 0,000 ... 45,000 & \centering  & \centering \arraybackslash 0,001 \\
\hline
\raggedright {'Description': ''} & \centering PereklyuchatelVkh2boolconst1S & \centering Выведено \\ Введено & \centering --- & \centering \arraybackslash --- \\
\hline
\raggedright {'Description': ''} & \centering Programmnyjklyuch1s & \centering Введено \\ Выведено & \centering --- & \centering \arraybackslash --- \\
\hline
\raggedright {'Description': ''} & \centering NezavisimoeRVnasrabatyvanie2Tvremya & \centering 0 ... 100000 & \centering  & \centering \arraybackslash 1 \\
\hline
\raggedright {'Description': ''} & \centering PereklyuchatelVkh4boolconst1S & \centering Откл. \\ БК-I-б \\ БК-I-м \\ БК-Z & \centering --- & \centering \arraybackslash --- \\
\hline
\raggedright {'Description': ''} & \centering NezavisimoeRVnasrabatyvanie1Tvremya & \centering 0 ... 100000 & \centering  & \centering \arraybackslash 1 \\
\hline
\raggedright {'Description': ''} & \centering Programmnyjklyuch4s & \centering Введено \\ Выведено & \centering --- & \centering \arraybackslash --- \\
\hline
\raggedright {'Description': ''} & \centering PereklyuchatelVkh4boolconst2S & \centering Откл. \\ БК-I-б \\ БК-I-м \\ БК-Z & \centering --- & \centering \arraybackslash --- \\
\hline
\raggedright {'Description': ''} & \centering Programmnyjklyuch2s & \centering Введено \\ Выведено & \centering --- & \centering \arraybackslash --- \\
\hline
\raggedright {'Description': ''} & \centering Programmnyjklyuch3s & \centering Введено \\ Выведено & \centering --- & \centering \arraybackslash --- \\
\hline
%===m
\end{longtable}
\normalsize

\end{enumerate}
%%%%%%%%%%%%%%%%%%%%%%%%%%%%%%%%%%%%%%%%%%%%%%%%%%%%%%%%%%%%%%%%%%%%%% КОНЕЦ ПЕРВАЯ СТУПЕНЬ ФАЗА-ЗЕМЛЯ %%%%%%%%%%%%%%%%%%%%%%%%%%%%%%%%%%%%%%%%%%%%%%%%%%


%%%%%%%%%%%%%%%%%%%%%%%%%%%%%%%%%%%%%%%%%%%%%%%%%%%%%%%%%%%%%%%%%%%%%% 2,3 СТУПЕНИ ФАЗА-ЗЕМЛЯ %%%%%%%%%%%%%%%%%%%%%%%%%%%%%%%%%%%%%%%%%%%%%%%%%%
\item\label{sec:secfz} Остальные ступени ДЗ от однофазных КЗ ДЗ-ФЗ-2, ДЗ-ФЗ-3
\begin{enumerate}[label=\arabic{section}.\arabic{subsection}.\arabic{enumi}.\arabic*, labelsep=4pt, leftmargin=0em, itemindent=65pt, parsep=0pt]

\item
Логика работы остальных ступеней ДЗ от однофазных КЗ выполнена одинаковой. Ввод в работу ступеней осуществляется программными ключами \textbf{«ДЗ-ФЗ-2(3)»}. Действие ступеней осуществляется с регулируемыми независимыми выдержками времени \textbf{«Тср ДЗ-ФЗ-2(3)»}. Контроль от БК задается отдельными программными переключателями \textbf{«БК ДЗ-ФЗ-2(3)»}.

\item
Разрешающий сигнал от БК для ступеней подхватывается от срабатывания РС соответствующей ступени и удерживается независимо от истечения времени ввода БК. Возврат разрешающего сигнала происходит только после возврата РС соответствующей ступени. На рисунке \ref{pict_dz:DZ_FZ_2} приведена функциональная схема логики ДЗ-ФЗ-2, выполнение функциональной схемы логики ступени ДЗ-ФЗ-3 аналогичное.

\vspace{2mm}
\begin{figure}[H]
\centering
\includegraphics[width=1\textwidth,height=1\textheight,keepaspectratio]{\fbpath/051.0041 ДЗ/img/DZ_FZ_2.pdf}
\captionof{figure}{Функциональная схема логики ДЗ-ФЗ-2}\label{pict_dz:DZ_FZ_2}
\end{figure}

\item
Уставки ступеней ДЗ от однофазных КЗ (ДЗ-ФЗ-2, ДЗ-ФЗ-3) представлены в таблице \ref{desc_dz:tbl7}.

\small
\begin{longtable}{|>{\centering\arraybackslash}m{5.3cm}|>{\centering\arraybackslash}m{3.3cm}|>{\centering\arraybackslash}m{4.2cm}|>{\centering\arraybackslash}m{1.8cm}|>{\centering\arraybackslash}m{1cm}|}
% Первая страница: полный заголовок
\caption{Параметры для настройки ступени ДЗ-ФЗ-2(3)\hfill\vspace{-0.5\baselineskip}}\label{desc_dz:tbl7}\\
\hline
\rowcolor{gray!30}
Параметр & Условное обозначение на схеме & Значение/ Диапазон & Единица измерения & Шаг \\
\hline
\endfirsthead

% Промежуточные страницы: «Продолжение таблицы»
\caption*{\hspace{3pt}\emph{Продолжение таблицы \ref{desc_dz:tbl7}}\hfill\ }\vspace{-0.5\baselineskip} \\
\hline
\rowcolor{gray!30}
Параметр & Условное обозначение на схеме & Значение/ Диапазон & Единица измерения & Шаг \\
\hline
\endhead

%===m>ДЗ|ДЗ-ФЗ-2
\raggedright {'Description': ''} & \centering RelesoprotivleniyadlyaDZFFN1Rsr & \centering 0,000 ... 100,000 & \centering  & \centering \arraybackslash 0,001 \\
\hline
\raggedright {'Description': ''} & \centering RelesoprotivleniyadlyaDZFFN1Xsr & \centering 0,000 ... 100,000 & \centering  & \centering \arraybackslash 0,001 \\
\hline
\raggedright {'Description': ''} & \centering RelesoprotivleniyadlyaDZFFN1F1 & \centering 0,000 ... 359,000 & \centering  & \centering \arraybackslash 0,001 \\
\hline
\raggedright {'Description': ''} & \centering RelesoprotivleniyadlyaDZFFN1F2 & \centering 0,000 ... 359,000 & \centering  & \centering \arraybackslash 0,001 \\
\hline
\raggedright {'Description': ''} & \centering RelesoprotivleniyadlyaDZFFN1F3 & \centering 0,000 ... 359,000 & \centering  & \centering \arraybackslash 0,001 \\
\hline
\raggedright {'Description': ''} & \centering RelesoprotivleniyadlyaDZFFN1NapravlennostDZFF & \centering Прямо \\ Обратно & \centering --- & \centering \arraybackslash --- \\
\hline
\raggedright {'Description': ''} & \centering RelesoprotivleniyadlyaDZFFN1Rng & \centering 0,000 ... 100,000 & \centering  & \centering \arraybackslash 0,001 \\
\hline
\raggedright {'Description': ''} & \centering RelesoprotivleniyadlyaDZFFN1Fng & \centering 0,000 ... 359,000 & \centering  & \centering \arraybackslash 0,001 \\
\hline
\raggedright {'Description': ''} & \centering PereklyuchatelVkh2boolconst1S & \centering Выведено \\ Введено & \centering --- & \centering \arraybackslash --- \\
\hline
\raggedright {'Description': ''} & \centering NezavisimoeRVnasrabatyvanie1Tvremya & \centering 0 ... 100000 & \centering  & \centering \arraybackslash 1 \\
\hline
\raggedright {'Description': ''} & \centering PereklyuchatelVkh4boolconst1S & \centering Откл. \\ БК-I-б \\ БК-I-м \\ БК-Z & \centering --- & \centering \arraybackslash --- \\
\hline
\raggedright {'Description': ''} & \centering Programmnyjklyuch1s & \centering Введено \\ Выведено & \centering --- & \centering \arraybackslash --- \\
\hline
%===m
\end{longtable}
\normalsize

\end{enumerate}
%%%%%%%%%%%%%%%%%%%%%%%%%%%%%%%%%%%%%%%%%%%%%%%%%%%%%%%%%%%%%%%%%%%%%% КОНЕЦ 2,3 СТУПЕНИ ФАЗА-ЗЕМЛЯ %%%%%%%%%%%%%%%%%%%%%%%%%%%%%%%%%%%%%%%%%%%%%%%%%%


%%%%%%%%%%%%%%%%%%%%%%%%%%%%%%%%%%%%%%%%%%%%%%%%%%%%%%%%%%%%%%%%%%%%%% БНН %%%%%%%%%%%%%%%%%%%%%%%%%%%%%%%%%%%%%%%%%%%%%%%%%%
\item\label{sec:secbnn} Взаимодействие с блокировкой при неисправностях цепей напряжения (БНН)
\begin{enumerate}[label=\arabic{section}.\arabic{subsection}.\arabic{enumi}.\arabic*, labelsep=4pt, leftmargin=0em, itemindent=65pt, parsep=0pt]

\item
Для исключения ложного действия при возникновении неисправности в цепях напряжения ступени ДЗ могут быть заблокированы от БНН. Блокировка от БНН вводится отдельно для каждой ступени ДЗ программными ключами \textbf{«БНН ДЗ-ФФ(ФЗ)-N»}.

\end{enumerate}
%%%%%%%%%%%%%%%%%%%%%%%%%%%%%%%%%%%%%%%%%%%%%%%%%%%%%%%%%%%%%%%%%%%%%% КОНЕЦ БНН %%%%%%%%%%%%%%%%%%%%%%%%%%%%%%%%%%%%%%%%%%%%%%%%%%


%%%%%%%%%%%%%%%%%%%%%%%%%%%%%%%%%%%%%%%%%%%%%%%%%%%%%%%%%%%%%%%%%%%%%% БК %%%%%%%%%%%%%%%%%%%%%%%%%%%%%%%%%%%%%%%%%%%%%%%%%%
\item\label{sec:secbk} Взаимодействие с блокировкой при качаниях (БК)
\begin{enumerate}[label=\arabic{section}.\arabic{subsection}.\arabic{enumi}.\arabic*, labelsep=4pt, leftmargin=0em, itemindent=65pt, parsep=0pt]

\item
Дистанционная защита контролируется блокировкой при качаниях, предотвращающей ложное срабатывание ДЗ при качаниях и асинхронном ходе. В устройстве реализовано два способа блокировки:
\begin{enumerate}[label=\arabic*)]
    \item по приращению токов прямой и обратной последовательности;
    \item по скорости изменения величины сопротивления dZ/dt.
\end{enumerate}

\item
Блокировка при качаниях, реагирующая на приращение во времени векторов тока прямой и обратной последовательностей, имеет два измерительных органа: чувствительный и грубый.

\item
Контроль ступеней ДЗ от блокировки при качаниях вводится программными переключателями \textbf{«БК ДЗ-ФФ(ФЗ)-N»}, имеющими четыре положения:
\begin{enumerate}[label=\arabic*)]
\item \textit{Откл.} -- ступень работает без контроля от БК;
\item \textit{БК-I-б} -- ступень пускается от сигнала ввода быстродействующих ступеней;
\item \textit{БК-I-м} -- ступень пускается от сигнала ввода медленнодействующих ступеней;
\item \textit{БК-Z} -- ступень пускается от сигнала БК по скорости изменения величины сопротивления.
\end{enumerate}
\end{enumerate}
%%%%%%%%%%%%%%%%%%%%%%%%%%%%%%%%%%%%%%%%%%%%%%%%%%%%%%%%%%%%%%%%%%%%%% КОНЕЦ БК %%%%%%%%%%%%%%%%%%%%%%%%%%%%%%%%%%%%%%%%%%%%%%%%%%


%%%%%%%%%%%%%%%%%%%%%%%%%%%%%%%%%%%%%%%%%%%%%%%%%%%%%%%%%%%%%%%%%%%%%% АУ ДЗ %%%%%%%%%%%%%%%%%%%%%%%%%%%%%%%%%%%%%%%%%%%%%%%%%%
\item\label{sec:secau} Автоматическое ускорение дистанционной защиты (АУ ДЗ)
\begin{enumerate}[label=\arabic{section}.\arabic{subsection}.\arabic{enumi}.\arabic*, labelsep=4pt, leftmargin=0em, itemindent=65pt, parsep=0pt]

\item
В составе ДЗ предусмотрена логика формирования пуска автоматического ускорения ДЗ (АУ ДЗ), выходной сигнал которого действует в логику автоматического ускорения. 
\item
Формирование сигнала пуска АУ ДЗ можно задать от одной из ступеней ДЗ-ФФ и ДЗ-ФЗ программными переключателями \textbf{«Ступень АУ ДЗ-ФЗ»} и \textbf{«Ступень АУ ДЗ-ФФ»}. В зависимости от положения программного переключателя \textbf{«Вывод напр. ДЗ при АУ»} сигнал пуска АУ ДЗ формируется либо от ненаправленного РС с охватом начала координат – положение \textit{«Вкл.»}, либо от направленного РС – положение \textit{«Откл.»}.

\item
При установке измерительных ТН на линии (определяется программным переключателем \textbf{«Место ТН»}) предусматривается автоматический перевод ускоряемой ступени ДЗ на ненаправленное РС с охватом начала координат.

\item
При необходимости для автоматически ускоряемой ступени ДЗ можно ввести контроль от БК и БНН программными переключателями \textbf{«БК АУ ДЗ»} и \textbf{«БНН АУ ДЗ»}.

\item
Функциональная схема логики АУ ДЗ приведена на рисунке \ref{pict_dz:AU_DZ}.

\vspace{2mm}
\begin{figure}[H]
\centering
\includegraphics[width=0.95\textwidth,height=0.95\textheight,keepaspectratio]{\fbpath/051.0041 ДЗ/img/AU_DZ.pdf}
\captionof{figure}{Функциональная схема  пуска АУ ДЗ}\label{pict_dz:AU_DZ}
\end{figure}

\item
Уставки автоматического ускорения ДЗ (АУ ДЗ) представлены в таблице \ref{desc_dz:tbl8}.

\small
\begin{longtable}{|>{\centering\arraybackslash}m{5.3cm}|>{\centering\arraybackslash}m{3.3cm}|>{\centering\arraybackslash}m{4.2cm}|>{\centering\arraybackslash}m{1.8cm}|>{\centering\arraybackslash}m{1cm}|}
% Первая страница: полный заголовок
\caption{Параметры для настройки АУ ДЗ\hfill\vspace{-0.5\baselineskip}}\label{desc_dz:tbl8}\\
\hline
\rowcolor{gray!30}
Параметр & Условное обозначение на схеме & Значение/ Диапазон & Единица измерения & Шаг \\
\hline
\endfirsthead

% Промежуточные страницы: «Продолжение таблицы»
\caption*{\hspace{3pt}\emph{Продолжение таблицы \ref{desc_dz:tbl8}}\hfill\ }\vspace{-0.5\baselineskip} \\
\hline
\rowcolor{gray!30}
Параметр & Условное обозначение на схеме & Значение/ Диапазон & Единица измерения & Шаг \\
\hline
\endhead

%===m>ДЗ|Пуск АУ ДЗ
\raggedright {'Description': ''} & \centering Ступень АУ ДЗ-ФФ & \centering Откл. \\ 1 ступ. \\ 2 ступ. \\ 3 ступ. \\ 4 ступ. \\ 5 ступ. \\ 6 ступ. & \centering --- & \centering \arraybackslash --- \\
\hline
\raggedright {'Description': ''} & \centering Место ТН & \centering Шины \\ Линия & \centering --- & \centering \arraybackslash --- \\
\hline
\raggedright {'Description': ''} & \centering Ступень АУ ДЗ-ФЗ & \centering Откл. \\ 1 ступ. \\ 2 ступ. \\ 3 ступ. & \centering --- & \centering \arraybackslash --- \\
\hline
\raggedright {'Description': ''} & \centering Вывод напр. ДЗ при АУ & \centering Вкл. \\ Откл. & \centering --- & \centering \arraybackslash --- \\
\hline
\raggedright {'Description': ''} & \centering PereklyuchatelVkh4boolconst3S & \centering Откл. \\ БК-I-б \\ БК-I-м \\ БК-Z & \centering --- & \centering \arraybackslash --- \\
\hline
\raggedright {'Description': ''} & \centering БНН АУ ДЗ & \centering Введено \\ Выведено & \centering --- & \centering \arraybackslash --- \\
\hline
%===m
\end{longtable}
\normalsize

\end{enumerate}
%%%%%%%%%%%%%%%%%%%%%%%%%%%%%%%%%%%%%%%%%%%%%%%%%%%%%%%%%%%%%%%%%%%%%% КОНЕЦ АУ ДЗ %%%%%%%%%%%%%%%%%%%%%%%%%%%%%%%%%%%%%%%%%%%%%%%%%%


%%%%%%%%%%%%%%%%%%%%%%%%%%%%%%%%%%%%%%%%%%%%%%%%%%%%%%%%%%%%%%%%%%%%%% ОУ ДЗ %%%%%%%%%%%%%%%%%%%%%%%%%%%%%%%%%%%%%%%%%%%%%%%%%%
\item\label{sec:secou} Оперативное ускорение дистанционной защиты (ОУ ДЗ)
\begin{enumerate}[label=\arabic{section}.\arabic{subsection}.\arabic{enumi}.\arabic*, labelsep=4pt, leftmargin=0em, itemindent=65pt, parsep=0pt]

\item
В устройстве реализовано два режима оперативного ускорения ДЗ: с выдержкой времени (ОУ ДЗ) и без выдержки времени (ОУ-бв ДЗ). Ввод в работу ОУ ДЗ осуществляется от внешнего сигнала \textit{«Ввод ОУ ДЗ»}. Выбор оперативно ускоряемой ступени ДЗ выполняется с помощью программных переключателей \textbf{«Ступень ОУ ДЗ-ФФ»} и \textbf{«Ступень ОУ ДЗ-ФЗ»}. Время срабатывания ОУ ДЗ задается уставкой \textbf{«Тср ОУ ДЗ»}.

\item
Ввод в работу ОУ-бв ДЗ осуществляется от внешнего сигнала \textit{«Ввод ОУ-бв ДЗ»}. Выбор оперативно ускоряемой ступени ДЗ без выдержки времени выполняется с помощью программных переключателей \textbf{«Ступень ОУ-бв ДЗ-ФФ»} и \textbf{«Ступень ОУ-бв ДЗ-ФЗ»}. При необходимости для ОУ-бв ДЗ также можно задать время срабатывания уставкой \textbf{«Тср ОУ-бв ДЗ»}.

\item
Кроме того, предусмотрена возможность автоматического ввода ОУ ДЗ и ОУ-бв ДЗ программными ключами \textbf{«Ввод ОУ ДЗ при выводе осн. защиты»} и \textbf{«Ввод ОУ-бв ДЗ при выводе осн. защиты»} при приеме внешнего сигнала о выводе основной защиты. Введенное ускорение может быть оперативно выведено от внешнего сигнала \textit{«Авт. введенное ОУ выведено»}.

\item
Необходимо обратить внимание, что как оперативное ускорение, так и оперативное ускорение без выдержки времени осуществляются от сигналов пуска ступени ДЗ, то есть контроль от БК и БНН для обоих режимов определяется соответствующими настройками самой ступени.

\item
Функциональные схемы логики ОУ ДЗ и ОУ-бв ДЗ приведены на рисунках \ref{pict_dz:OU_DZ} и \ref{pict_dz:OU_bv_DZ} соответственно.

\vspace{2mm}
\begin{figure}[H]
\centering
\includegraphics[width=1\textwidth,height=1\textheight,keepaspectratio]{\fbpath/051.0041 ДЗ/img/OU_DZ.pdf}
\captionof{figure}{Функциональная схема логики ОУ ДЗ}\label{pict_dz:OU_DZ}
\end{figure}

\vspace{2mm}
\begin{figure}[H]
\centering
\includegraphics[width=1\textwidth,height=1\textheight,keepaspectratio]{\fbpath/051.0041 ДЗ/img/OU_bv_DZ.pdf}
\captionof{figure}{Функциональная схема логики ОУ-бв ДЗ}\label{pict_dz:OU_bv_DZ}
\end{figure}

\item
Уставки оперативного ускорения ДЗ (ОУ ДЗ) представлены в таблице \ref{desc_dz:tbl9}. Уставки оперативного ускорения ДЗ без выдержки времени (ОУ-бв ДЗ) аналогичны уставкам ОУ ДЗ.

\small
\begin{longtable}{|>{\centering\arraybackslash}m{5.3cm}|>{\centering\arraybackslash}m{3.3cm}|>{\centering\arraybackslash}m{4.2cm}|>{\centering\arraybackslash}m{1.8cm}|>{\centering\arraybackslash}m{1cm}|}
% Первая страница: полный заголовок
\caption{Параметры для настройки ОУ ДЗ\hfill\vspace{-0.5\baselineskip}}\label{desc_dz:tbl9}\\
\hline
\rowcolor{gray!30}
Параметр & Условное обозначение на схеме & Значение/ Диапазон & Единица измерения & Шаг \\
\hline
\endfirsthead

% Промежуточные страницы: «Продолжение таблицы»
\caption*{\hspace{3pt}\emph{Продолжение таблицы \ref{desc_dz:tbl9}}\hfill\ }\vspace{-0.5\baselineskip} \\
\hline
\rowcolor{gray!30}
Параметр & Условное обозначение на схеме & Значение/ Диапазон & Единица измерения & Шаг \\
\hline
\endhead

%===m>ДЗ|ОУ ДЗ
\raggedright {'Description': ''} & \centering Ступень ОУ ДЗ-ФФ & \centering Откл. \\ 1 ступ. \\ 2 ступ. \\ 3 ступ. \\ 4 ступ. \\ 5 ступ. \\ 6 ступ. & \centering --- & \centering \arraybackslash --- \\
\hline
\raggedright {'Description': ''} & \centering Тср ОУ ДЗ & \centering 0 ... 100000 & \centering  & \centering \arraybackslash 1 \\
\hline
\raggedright {'Description': ''} & \centering Ступень ОУ ДЗ-ФЗ & \centering Откл. \\ 1 ступ. \\ 2 ступ. \\ 3 ступ. & \centering --- & \centering \arraybackslash --- \\
\hline
\raggedright {'Description': ''} & \centering БСТО ОУ ДЗ & \centering Введено \\ Выведено & \centering --- & \centering \arraybackslash --- \\
\hline
\raggedright {'Description': ''} & \centering Ввод ОУ ДЗ при выводе осн.защиты & \centering Введено \\ Выведено & \centering --- & \centering \arraybackslash --- \\
\hline
%===m
\end{longtable}
\normalsize

\end{enumerate}
%%%%%%%%%%%%%%%%%%%%%%%%%%%%%%%%%%%%%%%%%%%%%%%%%%%%%%%%%%%%%%%%%%%%%% КОНЕЦ ОУ ДЗ %%%%%%%%%%%%%%%%%%%%%%%%%%%%%%%%%%%%%%%%%%%%%%%%%%


%%%%%%%%%%%%%%%%%%%%%%%%%%%%%%%%%%%%%%%%%%%%%%%%%%%%%%%%%%%%%%%%%%%%%% БСТО %%%%%%%%%%%%%%%%%%%%%%%%%%%%%%%%%%%%%%%%%%%%%%%%%%
\item\label{sec:secbsto} Взаимодействие с блокировкой при протекании сквозного тока (БСТО)
\begin{enumerate}[label=\arabic{section}.\arabic{subsection}.\arabic{enumi}.\arabic*, labelsep=4pt, leftmargin=0em, itemindent=65pt, parsep=0pt]

\item
Для предотвращения неправильной работы защиты предусмотрена блокировка при протекании через трансформаторы тока двух выключателей сквозного тока внешнего КЗ, вызывающего неодновременное насыщение трансформаторов. Блокировка вводится отдельными программными ключами для первых ступеней ДЗ-ФЗ-1 и ДЗ-ФФ-1 (\textbf{«БСТО ДЗ-ФФ-1»} и \textbf{«БСТО ДЗ-ФЗ-1»} соответственно), а также для обоих режимов оперативного ускорения (\textbf{«БСТО ОУ ДЗ»} и \textbf{«БСТО ОУ-бв ДЗ»}).

\end{enumerate}
%%%%%%%%%%%%%%%%%%%%%%%%%%%%%%%%%%%%%%%%%%%%%%%%%%%%%%%%%%%%%%%%%%%%%% КОНЕЦ БСТО %%%%%%%%%%%%%%%%%%%%%%%%%%%%%%%%%%%%%%%%%%%%%%%%%%


\end{enumerate}
